\chapter{Discussão e Conclusões}
\label{cap:discussao}

\section{Discussão}
\label{sec:discussao}

Os \ScanTotalCasos{} casos do painel entre 2010 e 2025 apresentam distribuição espacial estruturada e autocorrelação positiva da taxa anual em todos os anos, com valor máximo de Moran em 2024. A elevação das incidências no Sul a partir da segunda metade do período é compatível com a mudança da barreira climática descrita para a região por \citeonline{lee2021}, que também analisaram a relação entre risco de dengue, urbanização e conectividade no Brasil.

A comparação central do trabalho é a varredura espacial com e sem padronização indireta por variáveis térmicas, sanitárias e urbanas, como exercício exploratório. O modelo sem ajuste delimita zonas de excesso sob o limite de 16\% e 200~km; zonas que permanecem ou mudam após cada ajuste indicam heterogeneidade residual, sem demonstrar causalidade \cite{jung2009glm}. O risco acima do esperado segue a formulação de varredura espacial \cite{kulldorff1997spatial}.

Na comparação entre as oito formulações, a coleção passou de 129 zonas sem ajuste para 128 no ajuste por calor, 178 no esgoto adequado, 136 na densidade urbana, 158 na proporção urbana, 193 na água, 195 no lixo e 184 no índice composto. O primário deslocou-se de Adolfo (SP, $RR = 2{,}99$) para Queluzito (MG, $RR = 2{,}08$) sob o ajuste por calor, para Iacri (SP) sob esgoto, proporção urbana, água e lixo, e para Araújos (MG) sob densidade urbana e índice composto. Todos os primários ficaram no Centro-Sul, e a sobreposição municipal com o modelo sem ajuste variou de 56\% (calor) a 82\% (densidade urbana): o eixo de sobre-risco do Centro-Sul manteve-se nas diferentes covariáveis, ainda que a posição do polo primário varie entre formulações.

O WMAPE sintetiza quanto cada padronização aproxima o esperado do observado no conjunto analisado: os valores caíram de 75,1\% sem ajuste para 64,9\% sob calor, 70,5\% sob água, 71,9\% sob esgoto, 71,9\% sob o índice composto, 72,1\% sob lixo, 74,3\% sob densidade urbana e 74,4\% sob proporção urbana. O cálculo é interno à amostra, de modo que a leitura é descritiva: valores menores apontam qual esquema de estratificação melhor acompanha a variabilidade geográfica basal da dengue, sem indicar capacidade preditiva. Um ajuste que absorvesse todo o excesso aproximaria o esperado do observado e reduziria o WMAPE ao custo de esvaziar a detecção; esse limite não foi atingido, pois todas as formulações mantiveram coleções extensas de zonas residuais, entre 128 e 195 zonas, cobrindo de 1.679 a 2.207 municípios com risco relativo primário acima de 2,0. A leitura conjunta das duas medidas orienta a escolha: o ajuste por calor combina a maior redução do WMAPE (10,2 pontos percentuais) com a manutenção de 128 zonas, ao passo que os ajustes por água, esgoto, lixo e índice composto reduzem o erro em menor grau e preservam polos com risco relativo mais alto, entre 2,58 e 2,86.

A decomposição regional mostra que a redução do WMAPE não é uniforme no território. Sob o ajuste por calor, o indicador cai de forma acentuada no Centro-Oeste (53,4\% para 38,7\%), no Sudeste (57,0\% para 43,7\%) e no Sul (101,9\% para 80,9\%), permanece estável no Nordeste (122,2\%) e aumenta no Norte (183,0\% para 245,9\%). O ajuste por água, por sua vez, reduz o indicador no Norte (183,0\% para 109,8\%) e no Nordeste (122,2\% para 106,8\%), com oscilações pequenas nas demais regiões. Como o WMAPE nacional pondera os casos, ele é dominado pelas regiões de maior volume e pode ocultar contrastes regionais como esses.

Sob o ajuste por calor, o polo primário deslocou-se de Adolfo (SP), no modelo sem ajuste, para Queluzito (MG). O risco relativo reportado foi de 2,99 na zona de Adolfo e de 2,08 na zona de Queluzito, que reuniu 362 municípios, ante 319 na zona primária sem ajuste. Como são zonas distintas, essa diferença não mede uma redução do risco no mesmo conjunto de municípios. A configuração das zonas mudou com a padronização térmica, mas essa comparação não estima quanto do excesso se deve ao calor. O aumento do WMAPE no Norte indica que o esquema acompanha menos de perto a distribuição observada nessa região, onde a incidência é mais baixa e a exposição ao calor é alta.

Zonas como Goiânia ($RR = 2{,}40$) e focos no interior de São Paulo, no Paraná e em Santa Catarina sob o modelo ajustado indicam heterogeneidade residual \cite{lee2021}. Fatores sociodemográficos, conectividade e saneamento são hipóteses plausíveis \cite{power2022,cardim2025}, mas não foram estimados neste estudo. A persistência territorial de dez ou mais anos acima de 300 por 100 mil não implica transmissão contínua na mesma população, pois mobilidade e imunidade podem alterar a composição dos grupos suscetíveis \cite{santos2019persistencia}. Nos estratos de saneamento, as taxas acumuladas foram maiores nos níveis com maior cobertura; para densidade urbana, o maior valor ocorreu no tercil intermediário (Seção~\ref{subsec:resultados_estratificacao}). Essas diferenças descritivas podem refletir a composição urbana e populacional dos grupos e a subnotificação diferencial, mas não estimam efeitos do saneamento ou da densidade.

Ao longo das oito formulações, 1.655 municípios (29,7\%) permanecem em zonas significativas em seis ou mais delas, e 967 em todas, com concentração no Centro-Sul (Figura~\ref{fig:robustez}). A persistência nesse eixo sustenta a leitura de heterogeneidade residual estável no Sul, no Sudeste e no Centro-Oeste, que é onde as coleções mais se sobrepõem entre as formulações.

A sensibilidade à exclusão de 2024 reforça a dependência do período acumulado: sem esse ano, os casos do painel caem para 15.664.546 e o número de zonas aumenta na maioria das formulações, com o polo primário deslocado para o interior de São Paulo em quase todas. A sobreposição municipal entre os cenários ficou entre 51\% e 63\%, e o WMAPE nacional aumentou em todas as formulações, indicando que a distribuição de 2010--2023 e 2025 é menos acompanhada pelo esquema de padronização, baseado em população e estratos fixos, do que a série completa. A reexecução sem o ano epidêmico foi motivada também pela investigação etiológica no estado do Rio de Janeiro, que identificou circulação de Oropouche em amostras de casos suspeitos de dengue no período \cite{costa2025oropouche}; a leitura do ano como dengue exclusiva deve, portanto, ser cautelosa.

\section{Conclusões}
\label{cap:conclusoes}

A taxa anual de incidência apresentou autocorrelação espacial positiva em cada um dos 16 anos, com índices globais de Moran entre 0,215 e 0,613 e máximo em 2024. Na análise acumulada sem ajuste, a varredura identificou 129 zonas sem sobreposição (nível nominal de 5\%), que reuniram 64,1\% dos casos do painel em 34,3\% dos municípios; o polo primário foi Adolfo (SP), com 319 municípios e $RR = 2{,}99$.

Os sete ajustes identificaram de 128 a 195 zonas, com deslocamento do primário (Queluzito sob calor, $RR = 2{,}08$; Iacri sob esgoto, proporção urbana, água e lixo; Araújos sob densidade urbana e índice composto). O ajuste por calor combinou a maior redução do erro absoluto percentual ponderado (de 75,1\% para 64,9\%) com a manutenção de 128 zonas residuais; os demais ajustes reduziram o erro entre 0,6 e 4,6 pontos percentuais. Essas comparações são exploratórias.

A reexecução sem 2024 mudou o número de zonas em todas as formulações e a posição do polo primário em sete delas, com sobreposição municipal entre os cenários de 51\% a 63\% e erro absoluto percentual ponderado maior em todos os casos, indicando dependência do período acumulado. A persistência de zonas de alta taxa no Sudeste, no Centro-Oeste e no Sul após as estratificações indica heterogeneidade residual, cuja atribuição a vulnerabilidades estruturais exige análise específica.

Adicionalmente, a sazonalidade e a distribuição regional informam o calendário das ações de vigilância \cite{pescarini2022}.

\section{Limitações e perspectivas}

Há limitações operacionais. Em primeiro lugar, a unidade de agregação ecológica municipal impede capturar a heterogeneidade intraurbana e as ilhas de calor intramunicipais, uma vez que a resolução do ERA5-Land ($\approx 9$~km) atenua contrastes de microescala em grandes metrópoles \cite{rodrigues2018}. Além disso, $T_{\text{máx}}$ é o máximo de quatro instantâneos diários, não o máximo horário contínuo. Em segundo lugar, os dados do SINAN estão sujeitos a atrasos de inserção e subnotificação diferencial; estudos de soroprevalência mostram que notificações sintomáticas podem subestimar substancialmente o número real de infectados \cite{perisse2020soroprevalencia}. Como limites de cobertura, 1.229 casos notificados em Fernando de Noronha, 121 casos em 39 códigos sem correspondência na malha de 2010 e 1.873 casos em 16 registros sem código municipal ficaram fora do painel espacial. Assim, as zonas identificadas devem ser interpretadas como contrastes relativos de notificações. O reporte hierárquico sem correção de multiplicidade torna os resultados exploratórios. Além disso, a varredura acumulada é influenciada pelo ano epidêmico de 2024, que responde por 29,1\% dos casos do painel; a reexecução sem esse ano alterou o número de zonas e os polos primários em sete das oito formulações, como descrito nos Resultados, e parte da coleção corresponde a excessos marginais, cuja significância reflete o volume acumulado sob o modelo de Poisson. A contagem de robustez entre as formulações atribui o mesmo peso a cada uma delas e não pondera a intensidade das zonas, de modo que um município pode manter-se em várias formulações por excessos marginais. Por fim, a janela circular da varredura, ainda que contida pelo teto duplo de 16\% da população e 200~km de raio, impõe restrições geométricas às zonas detectadas; estatísticas flexíveis podem representar melhor formas irregulares, embora também tenham limitações computacionais \cite{tango2005flexible}. Futuros desdobramentos devem incorporar modelos ecológicos explícitos e janelas com formatos flexíveis baseadas, por exemplo, em grafos de transporte.

%%%%%%%%%%%%%%%%%%%%%%%%
%%%INÍCIO DOS APENDICES%
%%%%%%%%%%%%%%%%%%%%%%%%
